As Large Language Models (LLMs) are increasingly deployed as autonomous agents in social and economic simulations, understanding their behavioral range becomes critical. Current research often focuses on "point-estimate" evaluations, measuring how well an LLM replicates the most likely human response. However, we argue that LLMs are better conceptualized as "multiverse simulators" capable of generating a superset of human behavioral outcomes. We propose a \method framework to explore the outcome manifold of \gpt across various social conflict scenarios. Our experiments demonstrate that while LLMs exhibit a strong "alignment bias" toward collaborative and safe behavior---manifesting as mode collapse in generic personas---they can successfully traverse the "tails" of the human behavioral distribution when properly conditioned. Specifically, we find that more strategic personas exhibit significantly higher behavioral entropy (1.01) compared to generic human personas (0.14), which overwhelmingly favor social harmony. Furthermore, we show that these simulated outcomes are grounded in underlying psychological and cultural factors, accurately shifting behavior manifolds to match specific sociological constraints. Our findings establish a formal framework for using LLMs to predict not just what *will* happen, but the full range of what *could* happen in complex social systems.
